\originalsection{18.315：作业8（原卷 Fall 2005）}
本组收录原题2、4. 一般分拆的多米诺铺砌生成函数及生日概率问题在清单中保留原题号,
不编入图论正文.

\begin{exercise}\label{mit:18315-hw8-q2}
\textbf{18.315, 原卷 Fall 2005, 作业8, 题2.}
固定一棵有 $n$ 个顶点、根为 $R$ 的树 $t$.
令 $h(v)$ 为以 $v$ 为根的子树顶点数, 包括 $v$ 自己;
$h(R)=n$, 非根叶点的 $h$ 值为1.
设 $A(t)$ 为顶点到 $[n]$ 的双射标号, 满足每个真祖先的标号小于其后代,
因此根标号必须为1. 证明
\[
|A(t)|=\frac{n!}{\prod_{v\in V(t)}h(v)}.
\]
\end{exercise}
\sourceof{官方 HW8 第1页题2; 编者将原路包含关系明确改述为祖先条件; 编者独立解答}
\begin{solution}
对 $n$ 归纳. $n=1$ 时只有一个标号, 公式为 $1!/1=1$.
设根的子树为 $t_1,\ldots,t_s$, 大小为 $n_1,\ldots,n_s$,
其和为 $n-1$. 根必须取1; 把其余标签分配给这些有区别的子树,
有 $\binom{n-1}{n_1,\ldots,n_s}$ 种.
对子树中的任意标签集, 保序重标为 $[n_i]$ 不改变合法标号数.
由归纳假设,
\begin{align*}
|A(t)|
&=\frac{(n-1)!}{\prod_i n_i!}
\prod_i\frac{n_i!}{\prod_{v\in V(t_i)}h(v)}\\
&=\frac{(n-1)!}{\prod_{v\ne R}h(v)}
=\frac{n!}{\prod_vh(v)},
\end{align*}
最后一步用 $h(R)=n$.
这个递推也给出双向构造: 合法全标号唯一决定各子树的标签集和子树标号,
任意这样的数据拼合后都满足全部祖先条件.
\end{solution}

\begin{exercise}\label{mit:18315-hw8-q4}
\textbf{18.315, 原卷 Fall 2005, 作业8, 题4.}
固定一个 Young 图形 $\lambda$.
以该形状的全部标准 Young 表为顶点, 若两个表恰在两个格子中的填数不同, 就连一条边.
证明这个图连通. 对同一根树形状的全部递增标号给出并证明对应的结论.
\end{exercise}
\sourceof{官方 HW8 第1页题4; 编者独立解答, 同一偏序交换引理同时覆盖两个要求}
\begin{solution}
先证明一个一般事实: 任意有限偏序的线性扩展,
可通过交换次序中相邻且不可比的两个元素彼此转化.
设当前扩展为 $a_1,\ldots,a_n$, 目标为 $b_1,\ldots,b_n$.
$b_1$ 是偏序的极小元, 故当前排在它前面的任何 $x$ 都不满足 $x<b_1$;
也不满足 $b_1<x$, 否则当前扩展已经违反偏序.
这些元素都与 $b_1$ 不可比, 因此可把 $b_1$ 逐步向左交换到第一位.
交换相邻不可比元素不改变它们与其他元素的先后关系, 仍是线性扩展.
固定第一位, 在删去 $b_1$ 的偏序上重复, 有限步得到目标.

对 Young 图形, 顶点为格子, 定义偏序由“左格在右格前、上格在下格前”生成.
标准 Young 表是把 $1,\ldots,n$ 填入这些格子, 沿每行和每列严格递增;
依标签从小到大列出格子, 恰为该偏序的线性扩展.
上面的相邻元素交换即交换数值 $j,j+1$ 所在的两个格子,
交换后仍是标准表, 且恰改变两个格子, 是原图允许的一条边.
所以这些边已经足够连通全部表.

对根树形状, 用祖先关系作偏序.
递增标号正是它的线性扩展, 同样的交换只在不可比顶点间进行,
保持祖先条件且恰改变两个标号. 因而相应的标号图也连通.
上述算法至多作 $n(n-1)/2$ 次相邻交换;
这是构造路径的统一上界, 不宣称每对表之间都达到此距离.
\end{solution}
